% =====================================================================
%  config/estilo.tex  --  FORMATAÇÃO DO TRABALHO (NÃO É PRECISO EDITAR)
% ---------------------------------------------------------------------
%  Cada configuração abaixo indica de onde veio a regra:
%  "Manual" = Manual para Elaboração de Trabalhos Científicos, Ibmec, 2022.
% =====================================================================

% ---------------------------------------------------------------------
% Pacotes básicos
% ---------------------------------------------------------------------
\usepackage[T1]{fontenc}        % acentuação correta no PDF
\usepackage[utf8]{inputenc}     % arquivos .tex salvos em UTF-8
\usepackage{lastpage}           % total de folhas
\usepackage{indentfirst}        % recua também o 1º parágrafo de cada seção
\usepackage{color}
\usepackage{graphicx}           % figuras (\includegraphics)
\usepackage{microtype}          % melhora a distribuição do texto
\usepackage{amsmath,amssymb}    % equações
\usepackage{booktabs}           % linhas de tabela com boa espessura
\usepackage{multirow}
\usepackage{array}
\usepackage{longtable}          % tabelas que ocupam mais de uma página
\usepackage{enumitem}           % listas (alíneas) personalizadas
\usepackage{listings}           % trechos de código-fonte
\usepackage{pdfpages}           % incluir PDF (ficha catalográfica)
\usepackage{textcase}           % \MakeTextUppercase
\usepackage{xcolor}

% ---------------------------------------------------------------------
% Fonte: Arial 12 (Manual, 3.2) -> usamos Helvetica, equivalente
% métrico da Arial disponível em qualquer instalação LaTeX.
% ---------------------------------------------------------------------
\usepackage[scaled=0.92]{helvet}
\renewcommand{\familydefault}{\sfdefault}

% ---------------------------------------------------------------------
% Margens (Manual, 3.2): esquerda e superior 3 cm; direita e inferior 2 cm
% Em modo frente-e-verso (twoside) o LaTeX espelha automaticamente.
% ---------------------------------------------------------------------
\setlrmarginsandblock{3cm}{2cm}{*}
\setulmarginsandblock{3cm}{2cm}{*}
% Número da página a ~2 cm da borda superior (Manual, 3.6)
\setheadfoot{20pt}{1cm}
\setheaderspaces{*}{0.55cm}{*}
\checkandfixthelayout

% ---------------------------------------------------------------------
% Parágrafos e espaçamento (Manual, 3.2 e 3.3)
% ---------------------------------------------------------------------
\setlength{\parindent}{1.25cm}  % recuo de parágrafo: 1,25 cm
\setlength{\parskip}{0pt}       % sem linha em branco entre parágrafos
\OnehalfSpacing                 % entrelinhas 1,5
% Tamanho "reduzido" (Arial 10): citações longas, notas de rodapé,
% legendas, fontes de ilustrações, paginação (Manual, Quadro 2)
\renewcommand{\ABNTEXfontereduzida}{\footnotesize}

% ---------------------------------------------------------------------
% Títulos das seções (Manual, 2.1.13 e 3.5) -- todos em Arial 12
%   1 SEÇÃO PRIMÁRIA         -> negrito, caixa alta
%   1.1 SEÇÃO SECUNDÁRIA     -> sem negrito, caixa alta
%   1.1.1 Seção terciária    -> negrito, caixa baixa
%   1.1.1.1 Seção quaternária-> itálico, caixa baixa
%   1.1.1.1.1 Seção quinária -> normal, caixa baixa
% ---------------------------------------------------------------------
\renewcommand{\ABNTEXchapterfont}{\sffamily\bfseries}
\renewcommand{\ABNTEXchapterfontsize}{\normalsize}
\renewcommand{\ABNTEXpartfont}{\sffamily\bfseries}
\renewcommand{\ABNTEXpartfontsize}{\normalsize}
\renewcommand{\ABNTEXsectionfont}{\sffamily\mdseries}
\renewcommand{\ABNTEXsectionfontsize}{\normalsize}
\renewcommand{\ABNTEXsubsectionfont}{\sffamily\bfseries}
\renewcommand{\ABNTEXsubsectionfontsize}{\normalsize}
\renewcommand{\ABNTEXsubsubsectionfont}{\sffamily\mdseries\itshape}
\renewcommand{\ABNTEXsubsubsectionfontsize}{\normalsize}
\renewcommand{\ABNTEXsubsubsubsectionfont}{\sffamily\mdseries\upshape}
\renewcommand{\ABNTEXsubsubsubsectionfontsize}{\normalsize}
\chapterstyle{abnt}  % reaplica o estilo com as fontes acima

% Espaço de 1,5 entre título e texto (Manual, 3.3)
\setlength{\afterchapskip}{1.5\baselineskip}
\setbeforesecskip{1.5\baselineskip}
\setaftersecskip{1.5\baselineskip}
\setbeforesubsecskip{1.5\baselineskip}
\setaftersubsecskip{1.5\baselineskip}
\setbeforesubsubsecskip{1.5\baselineskip}
\setaftersubsubsecskip{1.5\baselineskip}
\setbeforeparaskip{1.5\baselineskip}
\setafterparaskip{1.5\baselineskip}

% Numeração até a seção quinária (Manual, 3.7)
\setsecnumdepth{paragraph}  % paragraph = subsubsubsection (quinária)
\settocdepth{paragraph}

% ---------------------------------------------------------------------
% Sumário com a mesma tipografia do texto (Manual, 2.1.13)
% ---------------------------------------------------------------------
\renewcommand{\cftchapterfont}{\bfseries}
\renewcommand{\cftchapterpagefont}{\bfseries}
\renewcommand{\cftsectionfont}{\mdseries}
\renewcommand{\cftsectionpagefont}{\mdseries}
\renewcommand{\cftsubsectionfont}{\bfseries}
\renewcommand{\cftsubsectionpagefont}{\bfseries}
\renewcommand{\cftsubsubsectionfont}{\itshape}
\renewcommand{\cftsubsubsectionpagefont}{\itshape}
\renewcommand{\cftparagraphfont}{\normalfont}
\renewcommand{\cftparagraphpagefont}{\normalfont}
\renewcommand{\cftpartfont}{\bfseries}
\renewcommand{\cftpartpagefont}{\bfseries}
\setlength{\cftbeforechapterskip}{0.5em}
% Seções secundárias também em CAIXA ALTA no sumário (igual ao texto)
\makeatletter
\let\ABNTEXoldlsection\l@section
\renewcommand*{\l@section}[2]{\ABNTEXoldlsection{\MakeTextUppercase{#1}}{#2}}
\makeatother

% Títulos das partes de apêndices e anexos (Manual, 2.3.3 e 2.3.4)
\renewcommand{\apendicesname}{APÊNDICES}
\renewcommand{\anexosname}{ANEXOS}

% ---------------------------------------------------------------------
% Paginação (Manual, 3.6): só o número, em algarismos arábicos, no
% canto superior direito, fonte 10, sem cabeçalho com nome do capítulo.
% ---------------------------------------------------------------------
\makeevenhead{abntheadings}{\ABNTEXfontereduzida\thepage}{}{}
\makeoddhead{abntheadings}{}{}{\ABNTEXfontereduzida\thepage}
\makeheadrule{abntheadings}{\textwidth}{0pt}
\makeevenhead{abntchapfirst}{\ABNTEXfontereduzida\thepage}{}{}
\makeoddhead{abntchapfirst}{}{}{\ABNTEXfontereduzida\thepage}

% ---------------------------------------------------------------------
% Ilustrações e tabelas (Manual, 3.11 e 3.12)
% Título ACIMA (\caption) e fonte ABAIXO (\fonte{...}), em fonte 10.
% ---------------------------------------------------------------------
\captionnamefont{\ABNTEXfontereduzida}
\captiontitlefont{\ABNTEXfontereduzida}
\captiondelim{ -- }
% Numeração contínua no trabalho todo: Figura 1, 2, 3... (e não 2.1, 2.2)
\counterwithout{figure}{chapter}
\counterwithout{table}{chapter}
\counterwithout{equation}{chapter}
\setfloatlocations{figure}{htbp}

% Ambiente "quadro" (conteúdo textual com bordas fechadas) e sua lista
\newcommand{\quadroname}{Quadro}
\newcommand{\listofquadrosname}{Lista de quadros}
\newfloat{quadro}{loq}{\quadroname}
\newlistof{listofquadros}{loq}{\listofquadrosname}
\newlistentry{quadro}{loq}{0}
\renewcommand{\cftquadroname}{\quadroname\space}
\renewcommand*{\cftquadroaftersnum}{\hfill\textendash\hfill}
\setfloatlocations{quadro}{htbp}

% Ambiente "codigo" (trechos de código-fonte) e sua lista
\newcommand{\codigoname}{Código}
\newcommand{\listofcodigosname}{Lista de códigos}
\newfloat{codigo}{loc}{\codigoname}
\newlistof{listofcodigos}{loc}{\listofcodigosname}
\newlistentry{codigo}{loc}{0}
\renewcommand{\cftcodigoname}{\codigoname\space}
\renewcommand*{\cftcodigoaftersnum}{\hfill\textendash\hfill}
\setfloatlocations{codigo}{htbp}
\setfloatlocations{table}{htbp}

% ---------------------------------------------------------------------
% Código-fonte (útil principalmente no relatório técnico)
% ---------------------------------------------------------------------
\lstset{
  basicstyle=\ttfamily\footnotesize,
  breaklines=true,
  frame=tb,
  numbers=left,
  numberstyle=\tiny,
  showstringspaces=false,
  tabsize=2,
  extendedchars=true,
  inputencoding=utf8,
  literate={á}{{\'a}}1 {é}{{\'e}}1 {í}{{\'i}}1 {ó}{{\'o}}1 {ú}{{\'u}}1
           {Á}{{\'A}}1 {É}{{\'E}}1 {Í}{{\'I}}1 {Ó}{{\'O}}1 {Ú}{{\'U}}1
           {ã}{{\~a}}1 {õ}{{\~o}}1 {Ã}{{\~A}}1 {Õ}{{\~O}}1
           {â}{{\^a}}1 {ê}{{\^e}}1 {ô}{{\^o}}1 {Â}{{\^A}}1 {Ê}{{\^E}}1 {Ô}{{\^O}}1
           {à}{{\`a}}1 {À}{{\`A}}1 {ç}{{\c c}}1 {Ç}{{\c C}}1
}

% ---------------------------------------------------------------------
% Listas (alíneas) no padrão ABNT: a) b) c)
% ---------------------------------------------------------------------
\setlist[enumerate,1]{label=\alph*), leftmargin=1.25cm, itemsep=0pt, topsep=4pt}
\setlist[itemize,1]{leftmargin=1.25cm, itemsep=0pt, topsep=4pt}

% ---------------------------------------------------------------------
% Citações e referências -- sistema AUTOR-DATA (Manual, 3.8.2)
% Títulos destacados em NEGRITO nas referências (Manual, 2.3.1.1 e)
% ---------------------------------------------------------------------
\usepackage[alf, abnt-emphasize=bf, abnt-etal-cite=3, abnt-etal-list=0, abnt-full-initials=yes]{abntex2cite}

% ---------------------------------------------------------------------
% Links internos clicáveis, sem cor (texto deve ser preto -- Manual, 3.1)
% ---------------------------------------------------------------------
\usepackage{hyperref}
\hypersetup{hidelinks, bookmarksdepth=4, pdfencoding=auto}
\usepackage{bookmark}

% ---------------------------------------------------------------------
% Comandos auxiliares usados pelo template
% ---------------------------------------------------------------------
% Natureza do trabalho, conforme o formato escolhido no main.tex
\newcommand{\imprimirnatureza}{%
  \ifrelatorio
    Relatório técnico de projeto aplicado apresentado ao curso de
    \imprimircurso{} do \imprimirinstituicao, como requisito parcial
    para a obtenção do grau de \imprimirgrau.
  \else
    Monografia apresentada ao curso de \imprimircurso{} do
    \imprimirinstituicao, como requisito parcial para a obtenção do
    grau de \imprimirgrau.
  \fi
}

% Linha de assinatura da banca
\newcommand{\membrobanca}[2]{%
  \vspace{1.6cm}
  {\centering
   \rule{0.75\textwidth}{0.4pt}\\[2pt]
   {\SingleSpacing #1\\ #2\par}}
}

% Caixa de "nota para o aluno" -- APAGUE essas caixas na versão final!
% Liga/desliga em main.tex: \mostrarorientacoestrue / \mostrarorientacoesfalse
\newcommand{\notaaluno}[1]{\ifmostrarorientacoes%
  \par\vspace{0.5\baselineskip}\noindent\fbox{\begin{minipage}{\dimexpr\textwidth-2\fboxsep-2\fboxrule}%
  \SingleSpacing\footnotesize\textbf{Orientação ao aluno:} #1\par\vspace{2pt}\end{minipage}}\par\vspace{0.5\baselineskip}\fi
}

% Título completo: "TÍTULO: subtítulo" (o subtítulo só aparece se existir)
\newcommand{\imprimirtitulocompleto}{%
  \MakeTextUppercase{\imprimirtitulo}%
  \if\relax\detokenize\expandafter{\imprimirsubtitulo}\relax\else
    : \imprimirsubtitulo
  \fi
}

% Linha do coorientador na folha de rosto (só aparece se houver)
\makeatletter
\newcommand{\imprimircoorientadorlinha}{%
  \abntex@ifnotempty{\imprimircoorientador}{%
    \\ \imprimircoorientadorRotulo~\imprimircoorientador.}}
\makeatother

% Nomes usados pelo \autoref (ex.: "Quadro 1", "Código 1")
\providecommand{\quadroautorefname}{Quadro}
\providecommand{\codigoautorefname}{Código}
